\documentclass[10pt]{article}
\usepackage{icmlko}

\icmlmeta{IP 파운데이션 월드모델}{481}{막히는 것은 롱테일이 아니었다}{과거를 못 파는 것은 이름 없는 작품이 아니라 말이 너무 많은 작품이다 --- 그리고 가드 하나가 136번 반대 방향으로 걸려 있었다}

\begin{document}

\twocolumn[
\icmlheader

\begin{abstract}
직전 노트가 세 게이트 중 하나를 통과시키고 하나를 무효로 만들었다. 프레시
티처가 그 안에서 \textbf{가드가 반대 방향으로 걸린 자리}를 찾았다 --- 조항 59
는 \emph{``빈 것을 부정으로 읽지 마라''} 인데, 코드는 그 반대편을 하고 있었다:
\textbf{`확인된 0' 을 `모른다' 로 강등}. 무결과 표지를 결과 컨테이너의
\textbf{부재}로 잡은 탓이다.

고치니 디시 판정 불가가 \textbf{139 $\to$ 3}(남은 셋은 전부 망 예외), 실물
1,200단계 판정 불가가 \textbf{372(31.0\%) $\to$ 236(19.67\%)} 다. 그리고
실패 회계에서 \texttt{e.code} 를 되찾으니 --- 뉴스는 \textbf{503}, 네이버는
\textbf{403} 이었다. \emph{``레이트 리밋''} 은 측정이 아니라 \textbf{추정}이었다.

그 다음 노크가 방향을 뒤집었다. 디시 검색은 \textbf{쪽 파라미터가 있고 날짜
파라미터가 없다}. 상한은 작품당 \textbf{3,000건}이고 서버가 그 너머를
클램프한다. 2019 이전 도달은 \textbf{0/5}. 그런데 여덟 작품의 끝쪽을 세어
보니 --- \textbf{막히는 것은 롱테일이 아니라 인기작이었다.} 언급이 많은 작품은
3,000건이 \textbf{21--29일}밖에 안 덮고, 언급이 드문 작품은 검색이 전량을
소진해 \textbf{몇 해}를 판다.
\end{abstract}
]

\section{가드가 반대로 걸려 있었다}

조항 59 는 \emph{``성공 신호는 작업의 성공이 아니다''} 이고, 그중 제일 위험한
형태로 \emph{``빈 것을 부정으로 읽는 것''} 을 지목한다. 직전 노트는 그것을
\textbf{측정 전에 배선}했다 --- `확인된 0' 과 `판정 불가' 를 갈랐고, 그것을
미덕으로 적었다.

그런데 갈라 놓은 방향이 틀렸다. `dcin()' 이 무결과를 \textbf{결과 컨테이너의
부재}로 잡았는데, 디시의 진짜 무결과 표지는 \emph{``검색 결과가 없''} 이라는
\textbf{문구}다. 그래서 정상 응답 136건이 \textbf{`모른다' 로 강등}됐다.

\begin{figure*}[t]
\centering
\includegraphics[width=\textwidth]{figs/fig1.pdf}
\caption{왼쪽: 가드를 바로 걸자 판정 불가가 무너진다. 가운데: 여덟 작품 중
다섯이 서버 상한(120쪽 = 3,000건)에 붙었다. 오른쪽: 검색이 주는 총량이
많을수록 1쪽이 덮는 과거가 짧다 --- \emph{말이 많을수록 못 판다.}}
\label{fig:1}
\end{figure*}

전제를 캐시로 추정하지 않고 쟀다. 직전 노트가 죽인 질의 68 중 \textbf{12개를
다시 요청}해 무결과 표지 \textbf{12/12}(부트 $[1.0, 1.0]$). 그리고
\textbf{캐시에 원본 HTML 이 없어} `캐시 재해석' 이 사유 문자열 재분류였다는
한계를 \textbf{사전등록에 먼저 적었다}.

\claimx{`확인된 0' 을 `모른다' 로 강등하는 것도 조항 59 의 같은 병이다 ---
방향만 반대다.}{무결과 표지를 문구로 바꿔도 판정 불가가 안 줄면 원인은 다른
데 있다 --- 실측은 $139 \to 3$ 이었다.}

\section{그리고 `레이트 리밋' 은 추정이었다}

실패 회계가 \texttt{type(e).\_\_name\_\_} 만 적고 \texttt{e.code} 를 버렸다.
되찾아 보니 뉴스는 \textbf{503}(서버가 죽은 것) · 네이버는 \textbf{403}(거부)
이다. \emph{`레이트 리밋' 이라는 말 자체가 한 번도 측정된 적이 없었다.}

\textbf{그리고 더 나쁜 것을 찾았다.} 네이버 블로그 검색이 우리 UA 에
\textbf{HTTP 200 · 72,547바이트 · \texttt{id="notfound"}} 를 준다 ---
\textbf{체인소맨조차 0건}이다.

\claimx{무결과 표지는 차단의 얼굴로도 온다 --- `확인된 0' 과 `차단당한 0' 이
같은 화면이다.}{양성 통제(널리 알려진 질의)가 결과를 내면 그 0 은 진짜 0
이다 --- 실측은 체인소맨 0건이었고, 그래서 이 사이클의 블로그 수치는
\textbf{판정 불가}다.}

\section{방향이 뒤집혔다}

노크의 물음은 \emph{``디시로 과거를 팔 수 있는가''} 였다. 답은 셋이다.

\begin{center}\small
\begin{tabular}{ll}
\toprule
쪽 파라미터 & \textbf{있다}(\texttt{/post/p/N/q/}) \\
날짜 파라미터 & \textbf{없다}(정렬·페이징·폼 어디에도) \\
서버 상한 & \textbf{작품당 3,000건}(120쪽 $\times$ 25) \\
2019 이전 도달 & \textbf{0/5} \\
\bottomrule
\end{tabular}
\end{center}

여기까지는 예상대로였다. 뒤집힌 것은 \textbf{누가 막히는가}다.

\begin{center}\small
\begin{tabular}{lrr}
\toprule
작품 & 검색 총량 상한 & 1쪽이 덮는 과거 \\
\midrule
드래곤볼 & 3,000 (상한) & \textbf{29일} \\
브레이커 & 3,000 (상한) & \textbf{21일} \\
케이온! & 3,000 (상한) & 71일 \\
아포칼립스의 고인물 & 1,225 & 62일 \\
리블루밍 & 375 & 139일 \\
회귀한 엑스트라\ldots & \textbf{25} & \textbf{1,337일} \\
옥탑방 소드마스터 & \textbf{25} & \textbf{1,698일} \\
\bottomrule
\end{tabular}
\end{center}

여덟 작품 중 \textbf{다섯이 상한에 붙었고 셋은 검색이 전량을 소진}했다.
소진한 쪽이 \textbf{몇 해}를 판다.

\claimx{과거를 못 파는 것은 이름 없는 작품이 아니라 말이 너무 많은 작품이다
--- 창 폭은 언급 밀도의 함수다.}{끝쪽이 상한에 안 붙은 작품이 오히려 얕게
닿으면 이 설명은 틀렸다 --- 실측은 25건짜리가 1,698일을 팠다.}

\textbf{색인 바닥 가설}(아직 안 쟀다): 두 사이클이 독립한 뽑기로 각각
\textbf{2019-06-26} 과 \textbf{2019-06-09} 를 최솟값으로 냈다. 검색 색인 자체에
바닥이 있을 수 있다 --- 검정 설계를 원장에 적었고 요청 열 건이면 갈린다.

\section{작품별 갤은 채택 안 했다}

전용 갤 존재율 \textbf{9/30 = 0.30} $[0.1333, 0.4667]$. 사전등록 하한
$0.15$ 에 \textbf{구간 하한이 $0.0167$ 모자란다} --- \emph{점추정은 넘었지만
구간이 안 넘었으므로 채택하지 않는다}(노트 133). 분모도 못박았다: \textbf{7,982
가 아니라 한글 별칭 보유 2,390} 이다.

\section{이 사이클이 만든 잘못된 0 하나}

갤 파서가 아이콘 문자를 안 걷어 정확일치 \textbf{0/30} 을 냈다. 눈으로 상세를
읽지 않았으면 \emph{``전용 갤은 없다''} 를 능력 카드에 올릴 뻔했다. 고친 뒤
\textbf{9/30}.

\textbf{이 사이클의 주제가 「확인된 0 을 잘못 만드는 병」인데, 그 안에서 같은
병을 한 번 더 만들었다.} 그리고 실행 가드를 부정 조건으로만 걸어 요청이
\textbf{125건}(예산 100)까지 갔다 --- 27건이 같은 URL 재요청이었다.
\textbf{규약 위반으로 표시했다.}

\section{남는 것}

\textbf{「매칭 자」가 확정됐다.} 유효 수확으로 세면 실물 작품당 중앙값
\textbf{0}(200 중 \textbf{136 이 0}) --- \texttt{native} 13 / \texttt{synonyms}
1 / \texttt{english} 0 / \texttt{romaji} 0. \textbf{직전 노트의 수확 22.5 는
대부분 사라진다.} 프레시 티처가 층별로 갈라 낸 결과와 \textbf{독립으로 같은
결론}이다.

그러므로 다음 걸음은 크롤도 별칭도 아니다 --- \textbf{``이 작품에 대한 텍스트가
실제로 존재하는가'' 를 재는 자}다.

\end{document}
